\cotefile{157-3}
\scanpage{177}
\mspage{126}
commutatif dans $\mathrm{Hot}_{W}$ (mais non dans \unsure{$\mathit{Cat}$}, sauf si $X_0=\varnothing$). Supposons que $W$ satisfait $W(7\,\mathrm{bis})$\inserted{(axiome du carré cocartésien, cf p.~75)}. Alors on a
\[
\left\{\begin{aligned}
\alpha_1\in W&\Longleftrightarrow i_2\in W\\
\alpha_2\in W&\Longleftrightarrow i_1\in W .
\end{aligned}\right.
\]
% slanted
\margin{\uline{NB} Pour l'implication $\deleted{\unsure{\alpha}}\Rightarrow$ on n'a besoin que de la forme \unsure{douce A|douce} de $W(7\,\mathrm{bis})$, (i.e. \unsure{(avec} les notations de \uline{XVI}) de $\mathrm{Loc}\,(6\,\mathrm{bis})$}

\deleted{\unsure{Comme}} \uline{Dém} Utilisant le lemme précédent au cas du \unsure{cas|cor.}, on voit que chacune des flèches $\alpha_1,\alpha_2,i_1,i_2$ est dans $W$, ssi la flèche correspondante $\alpha'_1,\alpha'_2,i'_{\unsure{1|2}}$ ou $i'_2$ l'est. Donc les équivalences ci-dessus équivalent à
\[
\begin{aligned}
\alpha'_1\in W&\Longleftrightarrow i'_2\in W\\
\alpha'_2\in W&\Longleftrightarrow i'_1\in W ,
\end{aligned}
\]
lesquelles sont vraies, en vertu de $W(7\,\mathrm{bis})$ \inserted{(formulation \unsure{duale|dure})} appliqué \deleted{\unsure{aux}} à $\mathscr X$ et aux deux \deleted{\unsure{ouverts}}\uline{\unsure{fermés|souscatég.|segments}} $\,{}_{1/}\mathscr X_{\deleted{\unsure{/1}}}$, $\,{}_{2/}\mathscr X_{\deleted{\unsure{/2}}}$ de $\mathscr X$.

\uline{Corollaire} : Considérons \uline{deux} diagrammes de type $\Phi^0$, et un hom. \inserted{$f=(f_0,f_1,f_2)$} de l'un dans l'autre (cf diagramme à traits pleins ci-contre) :
\diagram{\begin{tikzcd}[ampersand replacement=\&,column sep=small,row sep=small]
\& X_1 \arrow[dddd,"f_1" near end,crossing over] \arrow[rrd,dashed,"i_1"] \& \& \\
X_0 \arrow[ru,"\alpha_1"] \arrow[rrd,"\alpha_2"] \arrow[dddd,"f_0"'] \& \& \& \mathscr X \arrow[dddd,dashed,"f\deleted{\unsure{\mathscr X}}"]\\
\& \& X_2 \arrow[ru,dashed,"i_2"] \arrow[dddd,"f_2"'] \& \\
\& \& \& \\
\& X'_1 \arrow[rrd,dashed,"i'_1"] \& \& \\
X'_0 \arrow[ru,"\alpha'_1"] \arrow[rrd,"\alpha'_2"] \& \& \& \mathscr X'\\
\& \& X'_2 \arrow[ru,dashed,"i'_2"] \&
\end{tikzcd}}
d'où un hom
% archive sheet number: 177
\scanpage{178}
\mspage{127}
\diagram{\begin{tikzcd}[ampersand replacement=\&,column sep=small,row sep=small]
\& \& X_1 \& \& \& X'_1\\
f_{\mathscr X}=\int f:\quad\mathscr X=\int \& X_0 \arrow[ru]\arrow[rd] \& \& {\longrightarrow}\quad\mathscr X'=\int \& X'_0 \arrow[ru]\arrow[rd] \& \\
\& \& X_2 \& \& \& X'_2\rlap{,}
\end{tikzcd}}
s'insérant dans un \deleted{cube} \inserted{diagramme cubique} (cf diagramme précédent complété par des pointillés). Considérons les conditions suivantes
\[
\begin{aligned}
\text{a)}\quad & f_0\in W\\
\text{b)}\quad & f_1\text{ et }f_2\text{ sont }\in W\\
\text{c)}\quad & f_{\mathscr X}=\int(f_0,f_1,f_2)\quad\in W
\end{aligned}
\]
Supposons l'axiome $W(7)$ (\textquotedblleft Lemme des 5\textquotedblright, cf page 76) satisfait. Alors deux \deleted{\unsure{quel}} parmi les conditions $a,b,c$ impliquent la troisième :
\[
\left\{\begin{aligned}
a\mathbin{\&} b&\Longrightarrow c\\
b\mathbin{\&} c&\Longrightarrow a\\
c\mathbin{\&} a&\Longrightarrow b .
\end{aligned}\right.
\]

En effet, on a un hom du \deleted{\unsure{diagra}} \unsure{mme|même} cubique ci-dessus dans le diagramme similaire, où $\deleted{\unsure{\alpha}}X_1,X_2,X'_1,X'_2$ sont remplacés respectivement par les \deleted{\unsure{ouverts}}\uline{\unsure{fermés|souscatég.}} $\,{}_{1/}\mathscr X_{\deleted{\unsure{/1}}},{}_{2/}\mathscr X_{\deleted{\unsure{/2}}}$ de $\mathscr X$, et $\,{}_{1/}\mathscr X'_{\deleted{\unsure{/1}}},{}_{2/}\mathscr X'_{\deleted{\unsure{/2}}}$ de $\mathscr X'$, les deux carrés sup. et inférieur étant les inclusions d'\unsure{ouverts|éléments} de $\mathscr X$ resp. de $\mathscr X'$. Ainsi, les conditions $a),b),c)$ \inserted{sont} équivalentes respectivement aux conditions similaires pour le nouveau diagramme cubique
% archive sheet number: 178
\scanpage{179}
\mspage{128}
($a)$ et $c)$ n'ayant d'ailleurs pas changé et $b$ devenant $b')$ $f'_1:\mathscr X_{/1}\to\mathscr X'_{/1}$ et $f'_2:\mathscr X_{/2}\to\mathscr X'_{/2}$ sont dans $W$). Alors la conclusion \inserted{du cor.} est équivalente à $W(7)$ \inserted{(formulation \unsure{duale|dure})} dans le cas du nouveau diagramme cubique, \unsure{qed.|q.e.d.}

3.2. \uline{Comparaison de}
\diagram{\begin{tikzcd}[ampersand replacement=\&,row sep=tiny,column sep=small]
\&X_1\\
\int X_0\arrow[ru]\arrow[rd]\&\\
\&X_2
\end{tikzcd}}
\uline{avec une somme amalgamée ordinaire}. \uline{Carrés $W$-homotopiquement cocartésiens}.

La catégorie intégrale
\diagram{\begin{tikzcd}[ampersand replacement=\&,row sep=tiny,column sep=small]
\&X_1\\
\mathscr X=\int X_0\arrow[ru]\arrow[rd]\&\\
\&X_2
\end{tikzcd}}
est solution du problème universel suivant, \deleted{\unsure{2}} variante 2-catégorique de celui que définit la limite inductive du diagramme : il s'agit, pour $Y$ objet variable dans $\mathit{Cat}$, de considérer les systèmes de trois foncteurs $f_0,f_1,f_2$ et de deux hom. fonctoriels $\lambda_1,\lambda_2$ :
\diagram{\begin{tikzcd}[ampersand replacement=\&,column sep=small,row sep=small]
\&\&X_0\arrow[dl,"\alpha_1"']\arrow[dr,"\alpha_2"]\arrow[dd,"f_0", ""{name=C,near end}]\&\\
(\ast)\&X_1\arrow[dr,"f_1"',""{name=L,near end}]\&\&X_2\arrow[dl,"f_2", ""{name=R,near end}]\\
\&\&Y\&\\
\arrow[from=L,to=C,"\lambda_1",shorten <=1pt,shorten >=2pt]
\arrow[from=R,to=C,"\lambda_2"',shorten <=1pt,shorten >=2pt]
\end{tikzcd}}
\[
\begin{aligned}
\lambda_1 &: f_1\alpha_1\to f_0\\
\lambda_2 &: f_2\alpha_2\to f_0
\end{aligned}
\]
% archive sheet number: 179
\scanpage{180}
\mspage{129}
Se donner une flèche dans $\mathit{Cat}$
\[
\mathscr X\longrightarrow Y
\]
\textquotedblleft revient au $=$\textquotedblright\ que de se donner un système $(f_0,f_1,f_2,\lambda_1,\lambda_2)$ comme dessus, \inserted{et ceci} \textquotedblleft fonctoriellement en $Y$\textquotedblright. C'est là \deleted{\unsure{le|les}} \unsure{cas|un} particulier de la détermination générale des foncteurs d'\deleted{une intégrale} la catégorie fibrée \textquotedblleft \unsure{intégrée|intégrale}\textquotedblright\ d'un \inserted{\unsure{contrav.}} foncteur
\[
I^0\longrightarrow\mathit{Cat},\qquad
\begin{aligned}i&\longmapsto X_i\\u&\longmapsto u^*\end{aligned}
\]
\inserted{vers une cat. $Y$,} en termes d'un système de foncteurs
\[
X_i\xrightarrow{f_i}Y ,
\]
et d'homomorphismes fonctoriels
\[
f_i\circ u^*\xrightarrow{\lambda_u} f_j
\]
pour $u:i\to j$ d'où $\unsure{u'|u^*}:X_j\to X_i$, avec des conditions de transitivité pour des flèches composées, dont on n'a pas à tenir compte dans le cas actuel, vu qu'il n'y a pas de flèches composées à regarder dans \deleted{$\unsure{\Phi}$} $I=\Phi$ (sauf des composés \deleted{\unsure{avec}} avec flèches identiques, \deleted{\unsure{quitte à}} qu'on peut se dispenser de \uline{regarder} ici \ldots).
% archive sheet number: 180
\scanpage{181}
\mspage{\deleted{\unsure{129}}\unsure{130}}
Si en particulier on a un diagramme \uline{commutatif} dans $\mathit{Cat}$
\diagram{\begin{tikzcd}[ampersand replacement=\&,row sep=small,column sep=small]
\&\& X_0 \arrow[dl,"\alpha_1"']\arrow[dr,"\alpha_2"]\&\\
Q={} \& X_1\arrow[dr,-,"f_1"']\&\& X_2\arrow[dl,"f_2"]\\
\&\&Y\&{,}
\end{tikzcd}}
on peut le compléter par \deleted{\unsure{$f_0$}}.
\[
f_0\overset{\mathrm{def}}{=}f_1\alpha_1=f_2\alpha_2
\]
en un diagramme
\diagram{\begin{tikzcd}[ampersand replacement=\&,row sep=small]
\&X_0\arrow[dl,"\alpha_1"']\arrow[dr,"\alpha_2"]\arrow[dd,"f_0", ""{name=C,near end}]\&\\
X_1\arrow[dr,"f_1"',""{name=L,near end}]\&\& X_2\arrow[dl,-,"f_2", ""{name=R,near end}]\\
\&Y\&{,}\\
\arrow[from=L,to=C,bend left=12,"\mathrm{id}"]
\arrow[from=R,to=C,bend right=12,"\mathrm{id}"']
\end{tikzcd}}
ce qui montre qu'on a une donnée \inserted{du type \unsure{$(\ast)$|cas|carré} ci-dessus,} définissant \inserted{donc} une flèche de
\diagram{\begin{tikzcd}[ampersand replacement=\&,row sep=tiny,column sep=small]
\& X_1\&\\
\unsure{\varphi|\rho}_Q:\quad\mathscr X=\int X_0\arrow[ru]\arrow[rd]\&{}\arrow[r]\&Y\\
\& X_2\&
\end{tikzcd}}
qui donne lieu à un diagramme cubique
% archive sheet number: 181
\scanpage{182}
\mspage{131}
\diagram{\begin{tikzcd}[ampersand replacement=\&,column sep=small,row sep=tiny]
\&X_1\arrow[rrd,"i_1", ""{name=T,near end}]\arrow[dddd,-,"\mathrm{id}" near end,crossing over]\&\&\\
X_0\arrow[ru,"\alpha_1"]\arrow[rrd,"\alpha_2"]\arrow[rrr,dashed,-,"i_0", ""{name=M,near end}]\arrow[dddd,-,"\mathrm{id}"']\&\&\&\mathscr X\arrow[dddd,"\unsure{\varphi|\rho}"]\\
\&\&X_2\arrow[ru,"i_2"', ""{name=B,near end}]\arrow[dddd,"\mathrm{id}"']\&\\
\&\&\&\\
\&X_1\arrow[rrd,"f_1"]\&\&\\
X_0\arrow[ru,"\alpha_1"]\arrow[rrd,"\alpha_2"]\arrow[rrr,dashed,"f_0"']\&\&\&Y\\
\&\&X_2\arrow[ru,"f_2"']\&{,}\\
\arrow[from=T,to=M,"\lambda_1"']\arrow[from=B,to=M,"\lambda_2"']
\end{tikzcd}}
commutatif sauf le carré supérieur, i.e. on a
\[
\varphi i_{\unsure{1|p}}=f_1,\quad \varphi i_2=f_2,\quad\text{et de plus }\varphi \deleted{\unsure{\alpha}}i_0=f_0,
\]
$\varphi$ transformant de plus les hom fonctoriels
\[
\lambda_1:i_1\alpha_1\to i_0,\qquad\lambda_2:i_2\alpha_2\to i_0
\]
en les \unsure{hom.} \deleted{flèches} \inserted{fonctoriels} identiques pertinents.
% slanted
\margin{Introduire $\varphi_Q$ dans une proposition en forme, et donner ses propriétés essentielles : flèches str. \unsure{de catégories|canoniques} sont transformées en les \unsure{fl.|$f_i$}, \unsure{homs|les} $\lambda_1$ et $\lambda_2$ \unsure{par|en} \uline{identiques}, $+$ les lemmes des p.~152, 154.}

\uline{Définition 3} : On dit que le diagramme commutatif
\diagram{\begin{tikzcd}[ampersand replacement=\&,column sep=small,row sep=tiny]
\&\&X_0\arrow[dl,"\alpha_1"']\arrow[dr,"\alpha_2"]\&\\
(\ast)\&X_1\arrow[dr,"f_1"']\&\&X_2\arrow[dl,"f_2"]\\
\&\&Y\&
\end{tikzcd}}
est $W$-homotopiquement cocartésien$^{\ast}$ \inserted{ou simplement $W$-cocartésien} \inserted{(ou qu'il définit $Y$ comme somme amalgamée $W$-homotopique des \unsure{etc.|$X_i$.})} si \unsure{l'hom|le hom} canonique
\diagram{\begin{tikzcd}[ampersand replacement=\&,row sep=tiny,column sep=small]
\&X_1\&\\
\mathscr X=\int X_0\arrow[ru]\arrow[rd]\&{}\arrow[r,"\varphi"]\&Y\\
\&X_2\&
\end{tikzcd}}
qu'il définit est dans $W$.
% archive sheet number: 182
\scanpage{183}
\mspage{132}
On a alors le

\uline{Corollaire 1} \inserted{Supposons $W(7\,\mathrm{bis})$.} Si le \deleted{diagramme} \inserted{carré} $(\ast)$ est \inserted{$W$-}homotopiquement cocartésien, alors
\[
\left\{\begin{aligned}
\alpha_1\in W&\Longleftrightarrow f_2\in W\\
\alpha_2\in W&\Longleftrightarrow f_1\in W .
\end{aligned}\right.
\]
Cela résulte aussitôt du diagramme cubique de la page précédente, \deleted{\unsure{où}} \deleted{\unsure{toutes}} les flèches verticales sont dans $W$, et de l'énoncé similaire concernant la face supérieure du cube, donné dans la prop.~13 p.~125.

\uline{Corollaire 2} \inserted{Supposons $W(7)$.} Considérons un diagramme cubique
\diagram{\begin{tikzcd}[ampersand replacement=\&,column sep=small,row sep=small]
X_0\arrow[r,"\alpha_1"]\arrow[rrd,"\alpha_2"]\arrow[ddd,-,"\deleted{f}\unsure{p|\rho|\varphi}_0"']\&X_1\arrow[rrd,"f_1"]\arrow[ddd,"\deleted{f}p_1",crossing over]\&\&\\
\&\&X_2\arrow[r,-,"f_2"]\arrow[ddd,"\deleted{f}p_2",crossing over]\&Y\arrow[ddd,"\deleted{f}p"]\\
\&\&\&\\
X'_0\arrow[r,-,"\alpha'_1"]\arrow[rrd,-,"\alpha'_2"]\&X'_1\arrow[rrd,"f'_1"]\&\&\\
\&\&X'_2\arrow[r,"f'_2"']\&Y'\rlap{,}
\end{tikzcd}}
\deleted{\unsure{où}} où les deux faces inf. et sup. du cube sont des carrés $W$-homotopiquement cocartésiens. Considérons les conditions
% archive sheet number: 183
\scanpage{184}
\mspage{133}
\[
\begin{aligned}
\text{a)}\quad &\unsure{p|\rho}_0\in W\\
\text{b)}\quad &p_1\text{ et }p_2\in W\\
\text{c)}\quad &p\in W
\end{aligned}
\]
Alors deux de ces conditions impliquent la troisième.

Résulte aussitôt de la \unsure{définition 3|définition 1}, et du résultat similaire (cor. page 126) pour le diagramme cubique similaire, avec $Y,Y'$ remplacés par
\diagram{\begin{tikzcd}[ampersand replacement=\&,row sep=tiny,column sep=small]
\&X_1\\
\mathscr X=\int X_0\arrow[ru]\arrow[rd]\&\\
\&X_2
\end{tikzcd}}
et
\diagram{\begin{tikzcd}[ampersand replacement=\&,row sep=tiny,column sep=small]
\&X'_1\\
\mathscr X'=\unsure{\cdot}\int X'_0\arrow[ru]\arrow[rd]\&\\
\&X'_2 .
\end{tikzcd}}

\unsure{3.3.} \uline{Intégration propre versus intégration lisse}. \uline{Énoncé d'équivalence entre les deux pts de vue}.
% slanted
\margin{$+$ \uline{commentaires sur le \deleted{\unsure{cas}} des limites $W$-homotopiques quelconques}.}

La définition 3, et les corollaires précédents, n'ont d'intérêt que dans la mesure où on dispose de critères maniables, permettant d'affirmer qu'un carré commutatif est bel et bien homotopiquement cocartésien. C'est ce qu'on va essayer de faire tantôt. Il y a aussi une perplexité, concernant le rôle privilégié joué, dans la théorie développée ici, par les immersions \deleted{\unsure{ouvertes}}\unsure{fermées}, à l'exclusion des immersions \deleted{fermées}\inserted{ouvertes}. Nous avons en somme \deleted{\unsure{recon}} reconnu, des
% archive sheet number: 184
\scanpage{185}
\mspage{\deleted{\unsure{133}}\unsure{134}}
point de vue homotopique, les diagrammes quelconques
\diagram{\begin{tikzcd}[ampersand replacement=\&,row sep=tiny,column sep=small]
\&X_1\\
X_0\arrow[ru]\arrow[rd]\&\\
\&X_2
\end{tikzcd}}
dans $\mathit{Cat}$, au cas \inserted{d'un diagramme} d'immersions \deleted{fermées}, où la somme amalgamée prend des allures familières. Mais pourquoi pas aussi bien les ramener au cas d'immersions \deleted{fermées}\inserted{ouvertes} ? On considérerait \unsure{donc|alors} le diagramme de départ comme \deleted{\unsure{formant}}\inserted{définissant} une catégorie \uline{cofibrée} sur
\diagram{\begin{tikzcd}[ampersand replacement=\&,row sep=tiny,column sep=small]
\&1\\
\Phi^0=\Psi=\quad 0\arrow[ru]\arrow[rd]\&\\
\&2
\end{tikzcd}}
soit $\deleted{\unsure{\mathscr Y}}\widetilde{\mathscr X}$. On a dans $\Psi$ les deux \deleted{fermés}\inserted{ouverts}
\[
\begin{aligned}
\Psi_{/1}&=U_1,&\Psi_{/2}&=U_2,&U_1\cap U_2&=\{0\},\\[-2pt]
\shortparallel&&\shortparallel\\[-2pt]
\{0,1\}&&\{0,2\}
\end{aligned}
\]
d'où dans $\widetilde{\mathscr X}$ deux ouverts images inverses
\[
\widetilde X_{/1},\ \widetilde X_{/2},\quad\text{avec}\quad \widetilde X_{/1}\cap\widetilde X_{/2}=X_0 ,
\]
d'où un diagramme \inserted{\unsure{comm.}} d'inclusions ouvertes
\diagram{\begin{tikzcd}[ampersand replacement=\&,row sep=tiny,column sep=small]
\&\&X_0\arrow[dl,"\alpha'_1"']\arrow[dr,"\alpha'_2"]\&\\
(\unsure{\ast|\ast\ast})\&\widetilde X_{/1}\arrow[dr,"i'_1"]\&\&\widetilde X_{/2}\arrow[dl,"i'_2"']\\
\&\&\widetilde{\mathscr X}\&{,}
\end{tikzcd}}
% archive sheet number: 185
\scanpage{186}
\mspage{135}
lequel reçoit le diagramme non comm.
\diagram{\begin{tikzcd}[ampersand replacement=\&,row sep=small]
\&X_0\arrow[dl,"\alpha_1"']\arrow[dr,"\alpha_2"]\arrow[dd,dashed,"i_0", ""{name=C,near end}]\&\\
X_1\arrow[dr,"i_1"', ""{name=L,near end}]\&\&X_2\arrow[dl,"i_2", ""{name=R,near end}]\\
\&\widetilde{\mathscr X}\&{,}\\
\arrow[from=C,to=L,dashed,"\deleted{\unsure{f_0|i_0}}\widetilde\lambda_1"{swap}]\arrow[from=C,to=R,dashed,"\widetilde\lambda_2"]
\end{tikzcd}}
les flèches
\[
X_1\to\widetilde{\mathscr X}_{/1},\qquad X_2\to\widetilde{\mathscr X}_{/2}
\]
étant dans $W$. (et $=$ $W$-coasphériques -- elles ont $=$ un adjoint à gauche). Donc on trouve, comme précédemment
\[
\left\{\begin{aligned}
\alpha_1\in W&\Longleftrightarrow i_2\in W\\
\alpha_2\in W&\Longleftrightarrow i_1\in W ,
\end{aligned}\right.
\]
et de $=$ on peut \deleted{développer} formuler un \textquotedblleft lemme du cinq\textquotedblright\ comme le cor. page 125, mais avec $\widetilde{\mathscr X}$, \inserted{$\widetilde{\mathscr X}'$} au lieu de $\mathscr X,\mathscr X'$. \deleted{\unsure{De $=$, $\widetilde{\mathscr X}$}}

De plus $\widetilde{\mathscr X}$ est la solution d'un problème universel \inserted{dans $\mathit{Cat}$}, en termes de \unsure{données}
\diagram{\begin{tikzcd}[ampersand replacement=\&,row sep=small]
\&X_0\arrow[dl,"\alpha_1"']\arrow[dr,"\alpha_2"]\arrow[dd,"f_0", ""{name=C,near end}]\&\\
X_1\arrow[dr,"f_1"', ""{name=L,near end}]\&\&X_2\arrow[dl,"f_2", ""{name=R,near end}]\\
\&Y\&{,}\\
\arrow[from=C,to=L,"\widetilde\lambda_1"']\arrow[from=C,to=R,"\widetilde\lambda_2"]
\end{tikzcd}}
% archive sheet number: 186
\scanpage{187}
\mspage{136}
la seule différence avec celui formulé à la p.~\deleted{\unsure{126|120}}\unsure{128|126} étant la direction opposée des flèches $\lambda_1,\lambda_2$. Pour \unsure{tt} carré commutatif
\diagram{\begin{tikzcd}[ampersand replacement=\&,row sep=small,column sep=small]
\&\&X_0\arrow[dl,"\alpha_1"']\arrow[dr,"\alpha_2"]\&\\
(\ast)\&X_1\arrow[dr,"f_1"']\&\&X_2\arrow[dl,-,"f_2"]\\
\&\&Y\&{,}
\end{tikzcd}}
on trouve encore une flèche canonique
\[
\widetilde{\unsure{\varphi|\rho}}:\widetilde{\mathscr X}\longrightarrow Y .
\]
Il serait \deleted{\unsure{tout}} aussi raisonnable de dire que le carré $(\ast)$ est $W$-homotopiquement cocartésien, lorsque $\widetilde\varphi$ est $\in W$, que d'adopter la définition 3, page 131 -- on trouverait exactement les $=\!{}^{\mathrm{s}}$ corollaires. Les deux définitions seraient-elles équivalentes ??

Je note que $\widetilde{\mathscr X}$ étant cofibrée sur $\Psi=\Phi^0$, $\widetilde{\mathscr X}^0$ est fibrée sur $\Psi^0=\Phi$. De façon précise, on aura
\diagram{\begin{tikzcd}[ampersand replacement=\&,row sep=tiny,column sep=small]
\&X_1^0\\
\widetilde{\mathscr X}^{0}=\int (X_{\unsure{0}})^0\arrow[ru,"\alpha_1^0"]\arrow[rd,"\alpha_2^0"']\&\\
\&X_2^0\rlap{,}
\end{tikzcd}}
% archive sheet number: 187
\scanpage{188}
\mspage{137}
\unsure{et} il s'agit donc de comparer les deux intégrales
\diagram{\begin{tikzcd}[ampersand replacement=\&,row sep=tiny,column sep=small]
\&X_1\&\&X_1^0\\
\mathscr X=\int X_0\arrow[ru,-,"\alpha_1"]\arrow[rd,"\alpha_2", "\deleted{\unsure{\alpha}}"']\&{,}\&\widetilde{\mathscr X}^0=\int X_0^0\arrow[ru,"\alpha_1^0"]\arrow[rd,"\alpha_2^0"]\&{,}\\
\&X_2\&\&X_2^0
\end{tikzcd}}
\unsure{notamment} comme objets de $\mathrm{HOT}_W(\Phi)$.

Pour cela, je reprends la construction \unsure{d'un|du} foncteur
\[
X\unsure{\longmapsto|\longleftarrow}\Phi(X):\quad\mathit{Cat}\longrightarrow\mathit{Cat}
\]
de I, p.~72\deleted{\unsure{,}}-75, mais je préfère ici noter ce foncteur $\Sigma(X)$ (car $\Phi$ désigne déjà une \deleted{petite} catégorie
\diagram{\begin{tikzcd}[ampersand replacement=\&,row sep=tiny,column sep=tiny]
\&1\\0\arrow[ru,-]\arrow[rd,-]\&\\\&2
\end{tikzcd}}
). On utilisera \unsure{l'exist. de|l'exist.|l'existence} homs fonctoriels en $X$ :
\diagram{\begin{tikzcd}[ampersand replacement=\&,row sep=small,column sep=small]
\&\Sigma(X)\arrow[dl,"p_X"']\arrow[dr,"q_X"]\&\\
X\&\&X^0
\end{tikzcd}}
qui sont $W$-asphériques, donc $\in W$. De plus, $\Sigma(X^0)=\Sigma(X)$, $p_{X^0}=q_X$, $q_{X^0}=p_X$.

Ainsi, \deleted{\illegible{1}} si
\diagram{\begin{tikzcd}[ampersand replacement=\&,row sep=tiny,column sep=small]
\& X_1\&\\
X_{\ast}=\Bigl(\,X_0\arrow[ru,"\alpha_1"]\arrow[rd,"\alpha_2"]\&\Bigr)\&\in\mathrm{Hom}(\Phi^0,\mathit{Cat})\\
\& X_2\&
\end{tikzcd}}
\diagram{\begin{tikzcd}[ampersand replacement=\&,row sep=tiny,column sep=small]
\&X_1^0\&\\
X_{\ast}^{0}=\Bigl(\,X_0^0\arrow[ru,-,"\alpha_1^0"]\arrow[rd,"\alpha_2^0"]\&\Bigr)\&\in\quad \unsure{\mathrm{id}|\mathrm{id}^{\mathrm{em}}}\rlap{,}\\
\&X_2^0\&
\end{tikzcd}}
% archive sheet number: 188
\scanpage{189}
\mspage{138}
on trouve
\diagram{\begin{tikzcd}[ampersand replacement=\&,row sep=small,column sep=small]
\&\Sigma(\unsure{\alpha_1|X_1})\&\\
\Sigma(X)_{\ast}=\Biggl(\ \Sigma(X_0)\arrow[ru,"\Sigma(\alpha_1)"]\arrow[rd,"\Sigma(\alpha_2)"]\&\&\Biggr)\\
\&\Sigma(X_2)\&
\end{tikzcd}}
et les flèches
\diagram{\begin{tikzcd}[ampersand replacement=\&,row sep=small,column sep=small]
\&\Sigma(X_\ast)\arrow[dl,"p_X"']\arrow[dr,"q_X"]\&\\
X_\ast\&\&X_\ast^0
\end{tikzcd}}
dans $\underline{\mathrm{Hom}}(\Phi^0,\mathit{Cat})$. Ces flèches sont dans $W$ \deleted{\unsure{donc}} terme à terme, donc \unsure{les} \deleted{\unsure{hom.}} hom dans $\mathit{Cat}$
\diagram{\begin{tikzcd}[ampersand replacement=\&,row sep=small,column sep=small]
\&\int\Sigma(X_{\unsure{\ast}})\arrow[dl,"\int p_X"']\arrow[dr,"\int q_X"]\&\\
\mathscr X=\int X_\ast\&\&\int X_\ast^0=\widetilde{\mathscr X}^0
\end{tikzcd}}
\deleted{\unsure{sont}} dans $W$.

Finalement, le \inserted{\unsure{\textquotedblleft innocent\textquotedblright|ci-devant|ci-dessus}} carré commutatif
\diagram{\begin{tikzcd}[ampersand replacement=\&,row sep=tiny,column sep=small]
\&\&X_0\arrow[dl]\arrow[dr]\&\\
(\ast)\quad Q\deleted{\unsure{\mathscr X}}\mathrel{=}\&X_1\arrow[dr]\&\&X_2\arrow[dl]\\
\&\&Y\&
\end{tikzcd}}
(et le diagramme correspondant $\widehat Q$
\diagram{\begin{tikzcd}[ampersand replacement=\&,row sep=tiny,column sep=small]
\&\&X_0\arrow[dl]\arrow[dr]\&\\
(\ast\ast)\&X_1\arrow[dr]\&\&X_2\arrow[dl]\\
\&\&\mathscr X=\int X_\ast\arrow[d,"\unsure{\varphi|\rho}_Q"]\&\\
\&\&Y\&
\end{tikzcd}}
donne naissance, en le transformant
% archive sheet number: 189
\scanpage{190}
\mspage{139}
par $X\mapsto X^0$ et par $X\mapsto\Sigma(X)$, et en tenant compte des hom fonctoriels en $X$
\diagram{\begin{tikzcd}[ampersand replacement=\&,row sep=tiny,column sep=small]
\&\Sigma(X)\arrow[dl]\arrow[dr]\&\\X\&\&X^0\rlap{,}
\end{tikzcd}}
à trois diagrammes du type $\ast\ast)$, i.e. modelés sur la catégorie opposée de
\diagram{\begin{tikzcd}[ampersand replacement=\&,row sep=small,column sep=small]
\&\&0\&\\
\Gamma={}\&1\arrow[ru,"\alpha_1"]\&\&2\arrow[lu,"\alpha_2"']\\
\&\&3\arrow[lu,"i_1"']\arrow[ru,"i_2"]\&\\
\&\&4\arrow[u,"\unsure{\varphi|\rho}"]\&
\end{tikzcd}}
% slanted
\margin{On \deleted{\unsure{impose}} a $\alpha_1i_1\ne\alpha_2i_2$, mais la relation $\alpha_1i_1\varphi=\alpha_2i_{\unsure{2|1}}\varphi$}
(i.e. \tikz[baseline=(a.base)]{\node[draw,ellipse,inner sep=2pt] (a) {$\Delta^1\times\Delta^1$};} auquel on a rajouté un objet initial),
% slanted
\margin{faux, car le carré n'est pas commutatif}
savoir
\[
\widehat Q,\qquad\widehat{\Sigma(Q)},\qquad\widehat{Q^0}
\]
et deux homomorphismes de diagrammes ($\in\underline{\mathrm{Hom}}(\Gamma,\mathit{Cat})$)
\diagram{\begin{tikzcd}[ampersand replacement=\&,column sep=small,row sep=small]
\&\&\widehat{\Sigma(Q)}\arrow[dl]\arrow[dr]\&\\
(\ast\ast\ast)\&\widehat Q\&\&\widehat{Q^0}
\end{tikzcd}}
(déduit de
\diagram{\begin{tikzcd}[ampersand replacement=\&,column sep=small,row sep=tiny]
\&\Sigma(Q)\arrow[dl]\arrow[dr]\&\\Q\&\&Q^0
\end{tikzcd}}
), lesquels sont des $W$ terme à terme (i.e. \deleted{dans} \unsure{dans} $W_\Gamma\subset\mathrm{Fl}(\underline{\mathrm{Hom}}(\Gamma,\mathit{Cat}))$). Explicitement, en les sommets $i=0,1,2$ \inserted{resp. $4$} de $\Gamma$, ceci n'est autre que
\diagram{\begin{tikzcd}[ampersand replacement=\&,column sep=tiny,row sep=tiny]
\&\&\Sigma(X_i)\arrow[dl,"p_{X_i}"']\arrow[dr,"q_{X_i}"]\&\&\&\&\Sigma Y\arrow[dl,"p_Y"']\arrow[dr,"q_Y"]\&\\
D_i=\Biggl(\&X_i\&\&X_i^0\Biggr)\&\text{resp.}\&Y\&\&Y^0
\end{tikzcd}}
en le sommet $3$ c'est
\diagram{\begin{tikzcd}[ampersand replacement=\&,column sep=small,row sep=small]
\&\int\Sigma X_\ast\arrow[dl,"\int p_{X_\ast}"']\arrow[dr,"\int q_{X_\ast}"]\&\\
\mathscr X=\int X_\ast\&\&\int X_\ast^0\simeq\widetilde{\mathscr X}^0
\end{tikzcd}}
obtenu par \textquotedblleft intégration\textquotedblright\ des \deleted{\unsure{flèch.}} flèches
\diagram{\begin{tikzcd}[ampersand replacement=\&,column sep=tiny,row sep=tiny]
\&\deleted{\unsure{\Sigma X_\ast}}\arrow[dl,"\deleted{\unsure{p_{X_\ast}}}"',"\backslash"{pos=.50,anchor=center}]\\
\deleted{\unsure{X_\ast}}\&
\end{tikzcd}}
\diagram{\begin{tikzcd}[ampersand replacement=\&,column sep=small,row sep=tiny]
\&\Sigma X_\ast\arrow[dl,"p_{X_\ast}"']\arrow[dr,"q_{X_\ast}"]\&\\
X_\ast\&\&X_\ast^0 .
\end{tikzcd}}
% archive sheet number: 190
\scanpage{191}
\mspage{140}
Regardant la restriction des deux morphismes de diagrammes $(\ast\ast\ast)$ aux parties des \inserted{trois} diagrammes de type $\Gamma^0$ au dessus de la queue $\{3,4\}$ de $\Gamma$, on trouve le diagramme à traits pleins que voici :
\diagram{\begin{tikzcd}[ampersand replacement=\&,row sep=small,column sep=small]
\&\int\Sigma X_\ast\arrow[dl,"W"']\arrow[dr,"W"]\arrow[dd,"{\varphi_{\Sigma Q}}"']\&\&\Sigma\widetilde{\mathscr X}\arrow[dl,"W"',postaction={decorate,decoration={snake,amplitude=.5pt,segment length=3pt}}]\arrow[dr,dashed,"W"]\arrow[dd,dashed,"{\Sigma\widetilde\varphi_Q}"]\&\\
\mathscr X\arrow[dd,"{\unsure{\varphi|\rho}_Q}"']\&\&\widetilde{\mathscr X}^0\arrow[dd,"{\varphi_{Q^0}=(\widetilde\varphi_Q)^0}"]\&\&\widetilde{\mathscr X}\arrow[dd,densely dotted,preaction={draw,solid,thin},"{\widetilde\varphi_Q}"]\\
\&\Sigma Y\arrow[dl,"W"']\arrow[dr,"W"]\&\&\Sigma Y\arrow[dl,dashed,"W"']\arrow[dr,dashed,"W"]\&\\
Y\&\&Y^0\&\&Y\rlap{,}
\end{tikzcd}}
que j'ai complété par un deuxième volet en pointillés. Dans ce diagramme de quatre carrés accolés, les côtés horizontaux des carrés \deleted{\unsure{représentent}} des flèches de $\mathit{Cat}$ qui sont \unsure{des|dans} $W$. On trouve ainsi :
\[
\varphi_Q\in W\Longleftrightarrow\varphi_{\Sigma Q}\in W\Longleftrightarrow
(\varphi_{Q^0}=\widetilde\varphi_Q^0)\in W\Longleftrightarrow
\Sigma\widetilde\varphi_Q\in W\Longleftrightarrow\widetilde\varphi_Q\in W .
\]
On trouve aussi que les cinq flèches verticales sont isomorphes, en tant que flèches de $\mathrm{Hot}_W$.

Pour résumer :

\uline{Proposition 14} Considérons un \inserted{diagramme} \unsure{(carré commutatif)} 
\diagram{\begin{tikzcd}[ampersand replacement=\&,column sep=small,row sep=tiny]
\&\&X_0\arrow[dl]\arrow[dr]\&\\
Q\deleted{\unsure{Q\to}}={}\&X_1\arrow[dr]\&\&X_2\arrow[dl]\\
\&\&Y\&
\end{tikzcd}}
dans $\mathit{Cat}$, et le diagramme \inserted{\textquotedblleft opposé\textquotedblright} déduit par $X\mapsto X^0$
\diagram{\begin{tikzcd}[ampersand replacement=\&,column sep=small,row sep=tiny]
\&\&X_0^0\arrow[dl]\arrow[dr]\&\\
Q^0={}\&X_1^0\arrow[dr]\&\&X_2^0\arrow[dl]\\
\&\&Y^0\&{,}
\end{tikzcd}}
% archive sheet number: 191
